%% kapa: exam=1 round=19 year=2008 period=1 form=A subject=appraisal_law subject_ko=부동산관계법규 seq=23 number=63
%% source: Q-Net 감정평가사 기출문제 (cst00309) artlSeq=1001962 file=extracted/1008/1001962/2008년 제19회 감정평가사 1차시험 기출문제/1교시.hwp
%% conversion: kordoc 4.24.51 -> tools/gen1.py (auto); review=proofread-sonnet; flags=none
\begin{question}{63}
부동산 가격공시 및 감정평가에 관한 법령상 ‘개별공시지가를 결정․공시하지 않아도 되는 토지’에 해당하지 않는 것은? (단, 관계법령에 의하여 지가의 산정 등에 개별공시지가를 적용하도록 규정되어 있는 토지와 시장․군수 또는 구청장이 관계행정기관의 장과 협의하여 개별공시지가를 결정․공시하기로 한 토지는 논외로 함)
\begin{choices}
  \item 표준지로 선정된 토지
  \item 국세 부과대상이 아닌 국유 공용토지
  \item 지방세 부과대상이 아닌 공유 공공용토지
  \item 농지전용부담금 부과대상이 아닌 토지
  \item 개발부담금 부과대상이 아닌 토지
\end{choices}
\end{question}
\answer{2}
